\section{Week in Review}
\sectionkicker{WEEKLY SYNTHESIS}
\begin{wideflow}
{\Large\bfseries 今週は、最前線が垂直化・検証・運用化へ動いた}\par
\smallskip
{\color{SurveyMuted}法律業務向けの構成、声で動く対話、検証つきの生物学利用、宣言・計測・提携を束ねる統治の弧、生成しない判断、コード調査の効率化、制作現場の画像が同時に前景化した。}\par
\end{wideflow}

\subsection*{今週何が変わったか}
\begin{enumerate}[leftmargin=1.6em]
\item \textbf{法律・音声・生物学が垂直方向に具体化した。} 9月17日の法律業務向け構成（2億3000万件超の検索基盤とValsの54.0\%というベンダー報告、Trusted Accessの統治つき）、9月15日の音声対話の新系列（97言語の行き来と背景実行、82.6や68.6\%といったベンダー側の測定値つき）、9月17日の検証つき生物学制度（標準と高リスクの二層、30日保持の事後監視つき）である。いずれも一次発表の範囲で読み、測定値の独立した再現には踏み込まない。
\item \textbf{宣言・計測・提携が一つの弧を描いた。} 9月の論考の三段の策、9月17日の研究所発信の三つの物差し（主導26\%、遮断率0.002\%、安全側6\%といった8月時点の切り取り）、9月18日のAccentureとの埋め込み評価の提携（各10億ドルの相互の見立て、非排他）である。論考が提携を生んだという因果ではなく、編集上の弧として読む。
\item \textbf{生成しない判断と調査の効率化と制作の画像が並んだ。} 9月15日のSystem OneとJev（文章生成を捨てた型つき判断、入力100万トークンあたり0.042ドル、型の範囲の保証つき）、9月16日のCode Scans（Agentic MapReduceの4段階、96\%や700時間超といった試用の値つき）、9月16日の2件の画像（Qwen-Image由来の一家とTsubaki.3の一般提供、創作道具と基盤モデルを競わせない読み分け）である。
\end{enumerate}

\subsection*{なぜそれらが一緒に重要なのか}
汎用の競い合いではなく、用途・検証・運用の具体化がそろって進んだ点にある。法律業務の構成、音声対話、生物学制度は、それぞれ異なる層の垂直化であり、単一のモデル性能の物語に束ねられない。宣言・計測・提携の弧は、速度の議論が運用の仕組みに降りてきたことを示す。非生成の判断、調査の効率化、制作の画像は、最前線の広がりが生成の外側と現場の双方に及んだことを示す。数値と利用条件、報告主体を切り離さない読みが条件になる。

この連鎖は製品間の直接的な互換性を意味しない。各公開、報道、一次記録は固有の主体・環境・測定条件を持つ。動作パラメーターの小ささを共通効率へ、利用条件の表示を共通ライセンスへ、速度やスコアを横断的な順位へ自動的に変換してはいけない。

\begin{communitynote}[Weekly Community Movement --- context only]
公開投稿では、法律・音声・生物学・計測、判断モデルへの関心が、公式発信に独立した言及が重なる形で観察される。これは今週の関心の広がりを示す背景であり、仕様や性能、提供条件、互換性を証明するものではない。締め切り後の速報2件は通常の集計には含めない。
\end{communitynote}

\clearpage
\subsection*{次に何を見るべきか}
次週以降は、この連鎖のどこがさらに具体化するかを追う。法律業務向けの価格や利用枠、提供地域、音声の対応言語や電話の範囲、LSVPの加入の広がりと高リスク枠の運用、埋め込み評価の他の評価者や他の開発者への広がり、研究所発信の計測の更新、Jevの再現や待ち行列の動き、Code Scansの試用の積み重ね、CAI-Image V2や動画化、Tsubaki Videoの10月20日までの特典後の動きはまだ詳細化されていない。さらに、ベンダーや研究機関の数値が、それぞれの条件から外へ一般化できるかを確認する必要がある。

\begin{claimboundary}[この総括の読み方]
各項目の公開条件、報告主体、測定環境は同じではない。ベンダーが示す速度、報道が伝える仕様、研究機関が報告する評価、話題で見られた関心は、それぞれ異なる種類の情報であり、単純な性能順位へまとめられない。未確定の点は、次の公開更新や再現可能な評価で確かめるべき観測点として残る。
\end{claimboundary}
